\input{_preamble}
%% Phantom bracket/paren/judgment alignment (opt-in).  When an object word
%% opens with delimiters or judgment marks, the gloss word below it is
%% padded by a \phantom of that leading run -- set in the OBJECT-tier font --
%% so its first real glyph sits under the object word's first real glyph.
%%
%% Test geometry.  Each probe column pairs an object word "<marks><STEM>"
%% with a gloss word "<STEM>" carrying the SAME stem glyphs.  When the
%% phantom is right, the gloss stem sits exactly under the object stem, so
%% their RIGHT edges coincide (pdftotext gives us that directly); when it is
%% off, the gloss stem is short of the object stem by the marks' width.
\begin{document}

% (1) alignment ON: gloss stem "Aaa" should end under object stem "Aaa".
\GlossPhantomAlign
\ex. \gll [Aaa BBB \\
     Aaa bbb \\
\glt `ON: single bracket.'

% (2) alignment ON, multiple leading marks "*([" before stem "Ddd".
\ex. \gll *([Ddd EEE \\
     Ddd eee \\
\glt `ON: three leading marks.'

% (3) alignment ON, gloss tier in \footnotesize: the shift must be the same
% full-size bracket width as in (1) (phantom is set in the object font).
% Scoped in a group so the smaller tier font cannot leak into (4)/(5).
{\GlossTierFont{2}{\footnotesize}%
\ex. \gll [Ccc FFF \\
     Ccc fff \\
\glt `ON: footnotesize gloss tier.'
}

% (4) alignment OFF: gloss stem "Bbb" sits at the column origin, under "[".
\GlossPhantomAlignOff
\ex. \gll [Bbb GGG \\
     Bbb ggg \\
\glt `OFF: gloss under the bracket.'

% (5) manual \GlossPhantom, automatic alignment still OFF: the manual
% phantom, set in the object-tier font, pushes gloss stem "Kkk" so its right
% edge lands under object stem "Kkk".  Brackets do not kern with letters, so
% the match is exact.
\ex. \gll [Kkk III \\
     \GlossPhantom{[}Kkk iii \\
\glt `Manual: gloss aligns past a bracket.'

% (6) auto ON must NOT fire on a macro-wrapped bracket: the scanner reads the
% control word (\textbf), not the "[" it hides, so no automatic phantom is
% added and the gloss stem stays at the column origin (left of the object
% stem by the bracket width).  Regression: the default character set once
% held stray backslashes (from writing \#\% in \str_set:Nn), which matched
% the leading control sequence of ANY macro-wrapped word and padded it.
\GlossPhantomAlign
\ex. \gll \textbf{[}Mmm NNN \\
     Mmm nnn \\
\glt `Auto does not see through a macro-wrapped bracket.'

%% (7)-(9) probe the SUBTRACTION: the pad is the object word's leading run
%% MINUS whatever leading run the tier's own word carries, clamped at zero.
%% Here the gloss word carries marks of its own, so object and gloss tokens
%% both start with them and the geometry is read off the TOKEN left edges
%% (the mark), not off the stem right edges as in (1)-(6): the stems differ
%% in case, and in (8) in size, so their widths are not comparable.

% (7) both tiers open with "(", both in the same font: the pad is zero and
% the two parens sit flush, which puts the stems flush too.  A translation
% reproducing the delimiters of what it translates is the ordinary case
% ("(e' questa)" glossed "(est celle-ci)"), and padding the gloss by the
% object's paren on top of its own put it one paren-width too far right --
% the aligner misaligning precisely the column that needed no help.
\ex. \gll (Ppp QQQ \\
     (ppp qqq \\
\glt `ON: same mark in both tiers, no pad.'

% (8) the same, with tier 2 at \tiny: the tier's own "[" is about half the
% object's, so the pad is the DIFFERENCE -- strictly between zero and the
% full bracket width of (1).  This is what separates the subtraction from
% the cheaper "skip the pad when the prefixes match": that rule would flush
% the two brackets and leave the real glyphs apart by the difference, which
% is the residual the phantom is set in the object font to remove.  \tiny
% rather than \footnotesize only so the gap clears TOL: at \footnotesize
% the difference is 0.37pt and no assertion could see it.
{\GlossTierFont{2}{\tiny}%
\ex. \gll [Rrr SSS \\
     [rrr sss \\
\glt `ON: same mark, smaller tier, pad is the difference.'
}

% (9) the clamp.  The gloss carries MORE mark than the object ("?" against
% "*(["), so aligning the real glyphs would need a negative pad, pulling the
% gloss's marks out of the column box and into the gap belonging to the
% column before it.  It is clamped at zero instead: the gloss stays at the
% column origin with its marks sticking out to the left of the object's.
\ex. \gll ?Ttt UUU \\
     *([ttt uuu \\
\glt `ON: pad clamped at zero, gloss keeps the column origin.'

\end{document}
